\exercisegroup

\begin{enumerate}
\def\labelenumi{\arabic{enumi})}
\item
  Let E be a complex normed space and f a symmetric bilinear form on the underlying real vector space \(E_{0}\) , such that \(f(x, x) = \|x\|^{2}\) for all \(x \in E\) . Show that there exists one and only one hermitian sesquilinear form g on E such that \(f(x, y) = \mathcal{R}g(x, y)\) (prove that \(f(x, iy) = -f(ix, y)\) by using formula (5) of V, p.~2), hence \(g(x, x) = \|x\|^{2}\) .
\item
  Let E be a real or complex normed space. Suppose that for every 2-dimensional vector subspace P (over R) in E, there exists a symmetric bilinear form \(f_{\mathbf{P}}\) defined on \(\mathbf{P} \times \mathbf{P}\) , such that \(f_{\mathbf{P}}(x,x) = \| x\|^2\) for all \(x \in \mathbf{P}\) . Show that \(f_{\mathbf{P}}\) is defined unambiguously and that there exists a hermitian sesquilinear form \(g\) on \(\mathbf{E} \times \mathbf{E}\) such that, for every real plane \(\mathbf{P} \subset \mathbf{E}\) , we have \(f_{\mathbf{P}}(x,y) = \mathcal{R}g(x,y)\) , hence \(g(x,x) = \| x\|^2\) . (If E is a real vector space, observe that we have \(\| x - y\|^2 + \| x + y\|^2 = 2(\| x\|^2 + \| y\|^2)\) for every pair of points of E, and deduce the identity
\end{enumerate}

\[
\left\| x + y + z \right\| ^ {2} - \left\| x + y \right\| ^ {2} - \left\| y + z \right\| ^ {2} - \left\| z + x \right\| ^ {2} + \left\| x \right\| ^ {2} + \left\| y \right\| ^ {2} + \left\| z \right\| ^ {2} = 0;
\]

if E is a complex vector space, apply exerc. 1.)

\begin{enumerate}
\def\labelenumi{\arabic{enumi})}
\setcounter{enumi}{2}
\tightlist
\item
  Let \(\mathbf{E}\) be a real finite dimensional vector space with dimension \(n, f\) a positive and separating symmetric bilinear form on \(\mathbf{E}\) , and \(\mathbf{B}_f\) the bounded convex set defined by the relation \(f(x, x) \leqslant 1\) . If \(\mathbf{a} = (a_1, \ldots, a_n)\) is a basis of \(\mathbf{E}\) and \(\Delta\) the discriminant of \(f\) with respect to this basis, we call the volume of \(\mathbf{B}_f\) with respect to \(\mathbf{a}\) the number \(v_{\mathbf{a}}(f) = \gamma_n |\Delta|^{-1/2}\) , where \(\gamma_n = \pi^{n/2} / \Gamma\left(\frac{n}{2} + 1\right)\) . If \(\mathbf{b} = (b_1, \ldots, b_n)\) is a second basis of \(\mathbf{E}\) , and if
\end{enumerate}

\[
a _ {1} \wedge a _ {2} \wedge \dots \wedge a _ {n} = \delta b _ {1} \wedge b _ {2} \wedge \dots \wedge b _ {n},
\]

we have \(v_{\mathbf{b}}(f) = |\delta| v_{\mathbf{a}}(f)\) .

\begin{enumerate}
\def\labelenumi{\alph{enumi})}
\item
  Show that, if \(f\) and \(g\) are two positive, separating symmetric bilinear forms such that \(\mathbf{B}_f \subset \mathbf{B}_g\) (which is equivalent to \(g \leqslant f\) ), then \(v_{\mathbf{a}}(f) \leqslant v_{\mathbf{a}}(g)\) (consider a basis for \(E\) which is orthogonal for both \(f\) and for \(g\) ).
\item
  Let A be a symmetric compact convex set in E, with 0 as an interior point. Show that among all positive, separating symmetric bilinear forms \(f\) on \(\mathbf{E}\) such that \(\mathbf{A} \subset \mathbf{B}_f\) , there exists one and only one for which the volume of \(\mathbf{B}_f\) (with respect to the given basis of \(\mathbf{E}\) ) is the smallest possible. (To show uniqueness observe that if \(\mathbf{A}\) is contained in \(\mathbf{B}_f\) and \(\mathbf{B}_g\) , it is in \(\mathbf{B}_{(f + g)/2}\) and that we have \(v_{\mathbf{a}}((f + g)/2) \leqslant \frac{1}{2}(v_{\mathbf{a}}(f) + v_{\mathbf{a}}(g))\) for every basis \(\mathbf{a}\) of \(\mathbf{E}\) which is orthogonal for both \(f\) and \(g\) .)
\item
  Let A be a symmetric compact convex set in E, with 0 as an interior point, and let \(f\) be the positive, separating, symmetric bilinear form such that \(\mathbf{A} \subset \mathbf{B}_f\) and that \(\mathbf{B}_f\) has the smallest possible volume with respect to a given basis of E. Show that there exist points \(x_1, \ldots, x_n, u_1, \ldots, u_n\) in E with the following properties:
\end{enumerate}

α) For every k, we have \(x_{k} \in A\) and \(f(x_{k}, x_{k}) = 1\) .

β) The sequence \((u_{1},\dots,u_{n})\) is an orthonormal basis of E for \(f\) .

\(\gamma)\) If we put \(x_{k} = \sum_{j=1}^{n} a_{kj} u_{j}\) for \(1 \leqslant k \leqslant n\) , we have \(a_{kj} = 0\) for \(k < j\) and \(a_{kk}^{2} \geqslant (n - k + 1)/n\) . (Argue by induction on \(k\) . Suppose \(x_{1}, \ldots, x_{k}, u_{1}, \ldots, u_{k}\) have been constructed, let \(P_{k}\) be the orthoprojector (for \(f\) ) from E onto the subspace generated by \(u_{1}, \ldots, u_{k}\) ; for every \(\varepsilon > 0\) , consider the bilinear form \(f_{\varepsilon}\) defined by

\[
f _ {\varepsilon} (x, y) = (1 + \varepsilon) ^ {k - n} f \left(P _ {k} x, P _ {k} y\right) + (1 + \varepsilon) ^ {k} f \left(x - P _ {k} x, y - P _ {k} y\right)
\]

and prove that \(\mathbf{A} \not\subset \mathbf{B}_{f_{\varepsilon}}\) using \(b\) ). For every integer \(p \geqslant 1\) choose a point \(y_{p}\) in \(\mathbf{A}\) not belonging to \(\mathbf{B}_{f_{1/p}}\) ; take for \(x_{k+1}\) a limit point of the sequence \((y_{p})\) such that

\[
k _ {f} \left(x _ {k + 1} - P _ {k} x _ {k + 1}, x _ {k + 1} - P _ {k} x _ {k + 1}\right) \geqslant (n - k) f \left(P _ {k} x _ {k + 1}, P _ {k} x _ {k + 1}\right);
\]

next choose \(u_{k+1}\) .

\(d\) ) Prove the analogues of \(b\) ) and \(c\) ) for the positive separating symmetric bilinear forms such that \(\mathbf{B}_f\subset \mathbf{A}\) and for which the volume of \(\mathbf{B}_f\) (with respect to a given bases of E) is the largest possible.

¶ 4) a) Let E be a real or complex normed space, of dimension ≥ 2, having the following property: the relation \textbar\textbar x\textbar\textbar{} = \textbar\textbar y\textbar\textbar{} implies the inequality

\[
\left\| x + y \right\| ^ {2} + \left\| x - y \right\| ^ {2} \leqslant 2 \left(\left\| x \right\| ^ {2} + \left\| y \right\| ^ {2}\right).
\]

Show that the norm on E is prehilbertian. (Reduce to the case when E is real and of dimension 2, by means of exerc. 1 and 2 of V, p.~60. In this case, let A be the unit ball of E, and let f be the positive, separating, symmetric bilinear form such that \(A \subset B_{f}\) and such that the volume of \(B_{f}\) with respect to a given basis is the smallest possible. Let \(x_{1}, x_{2}\) be the two points constructed in exerc. 3, c). Show that the points of intersection of the circle \(f(z, z) = 1\) and the bisectors of the two vectors \(x_{1}, x_{2}\) also belong to A, and conclude, by iteration, that \(A = B_{f}\) .)

\begin{enumerate}
\def\labelenumi{\alph{enumi})}
\setcounter{enumi}{1}
\item
  Let E be a real or complex normed space, of dimension \(\geqslant 2\) , having the following property: the relation \(\| x + y \| = \| x - y \|\) implies \(\| x + y \|^2 = \| x \|^2 + \| y \|^2\) . Show that the norm on E is prehilbertian (reduce to \(a\) )).
\item
  Prove the analogue of \(a\) when we assume that the relation \(\| x\| = \| y\|\) implies the inequality
\end{enumerate}

\[
\left\| x + y \right\| ^ {2} + \left\| x - y \right\| ^ {2} \geqslant 2 \left(\left\| x \right\| ^ {2} + \left\| y \right\| ^ {2}\right)
\]

(use exerc. 3, d)).

\begin{enumerate}
\def\labelenumi{\arabic{enumi})}
\setcounter{enumi}{4}
\tightlist
\item
  Let E be a real or complex vector space, of dimension \(\geqslant 2\) . Suppose a mapping \(x \mapsto \|x\|\) from E into \(R_{+}\) is given, satisfying : \(\|\lambda x\| = |\lambda|\) . \(\|x\|\) for every scalar \(\lambda\) , that \(\|x\| = 0\) implies x = 0 and that we have the « ptolemaic inequality »
\end{enumerate}

\[
\| a - c \| \cdot \| b - d \| \leqslant \| a - b \| \cdot \| c - d \| + \| b - c \| \cdot \| a - d \|
\]

for every \(a, b, c, d\) in E.

\begin{enumerate}
\def\labelenumi{\alph{enumi})}
\item
  Show that \(\| x\|\) is a norm on E (replace \(d\) by 0 and \(b\) by \(-a\) ).
\item
  Show that the norm on E is prehilbertian. (From the ptolemaic inequality deduce the inequality \(\|x+y\|^{2}+\|x-y\|^{2}\geqslant4\|x\|. \|y\|\) and use exerc. 4, c).) Prove the converse (show that in a hilbertian space, if we put \(a' = a/\|a\|^{2}\) , \(b' = b/\|b\|^{2}\) , we have the equality \(\|a'-b'\|=\|a-b\|/\|a\|. \|b\|\) ). If \(\|a\|=\|b\|=\|c\|=\|d\|\) and if the four vectors a, b, c, d are in the same plane, the two members of the ptolemaic inequality are equal.
\end{enumerate}

\begin{enumerate}
\def\labelenumi{\arabic{enumi})}
\setcounter{enumi}{5}
\tightlist
\item
  \begin{enumerate}
  \def\labelenumii{\alph{enumii})}
  \tightlist
  \item
    Let E be a real or complex normed space, of dimension \(\geqslant 2\) . Show that for every \(x \neq 0\) in E and every real number \(\alpha > 0\) , there exists an element \(y\) in E such that \(\| y \| = \alpha\) and \(\| x + y \|^2 = \| x \|^2 + \| y \|^2\) .
  \end{enumerate}
\end{enumerate}

\begin{enumerate}
\def\labelenumi{\alph{enumi})}
\setcounter{enumi}{1}
\item
  Suppose that if the vectors \(x, y\) in \(\mathbf{E}\) satisfy the relation \(\| x + y\|^2 = \| x\|^2 + \| y\|^2\) , then we also have \(\| x - y\|^2 = \| x\|^2 + \| y\|^2\) . Show that the norm on \(\mathbf{E}\) is prehilbertian. (Using a), reduce to the case in exerc. 4: restricting to the case where \(\mathbf{E}\) is 2-bidimensional, we prove that if \(\| x_1\| = \| x_2\| = 1\) , \(y = \frac{1}{2}(x_1 - x_2)\) and if \(z \in \mathbf{E}\) is such that \(\| y\|^2 + \| z\|^2 = \| y + z\|^2 = 1\) , then \(z = \frac{1}{2}(x_1 + x_2)\) or \(z = -\frac{1}{2}(x_1 + x_2)\) .)
\item
  Suppose that, for every vector \(x \neq 0\) in \(E\) , the set \(H\) of all vectors \(y\) satisfying \(\| x - y \|^2 = \| x \|^2 + \| y \|^2\) is stable under addition. Show that the conclusion of \(b\) ) holds. (Reduce to the case where \(E\) is real and of dimension 2. Using \(a\) ) and the compactness of the unit ball in \(E\) , show that \(H\) is a closed set containing at least two distinct half-lines with origin 0; prove that these two half lines are opposite to each other by observing that, if not, the convex set which they generate would be contained in \(H\) and will contain either \(x\) or \(-x\) .)
\end{enumerate}

\begin{enumerate}
\def\labelenumi{\arabic{enumi})}
\setcounter{enumi}{6}
\tightlist
\item
  Let E be a real or complex normed space, with dimension \(\geqslant 2\) , having the following property: there exists a real number \(\gamma\) distinct from 0 and from \(\pm 1\) , such that the relation \(\| x + y \| = \| x - y \|\) implies \(\| x + \gamma y \| = \| x - \gamma y \|\) .
\end{enumerate}

\begin{enumerate}
\def\labelenumi{\alph{enumi})}
\item
  Show that if \(\|x+y\|=\|x-y\|\) , the convex mapping \(\phi:\xi\mapsto\|x+\xi y\|\) from R into R is not constant on any interval. (Argue by reductio ad absurdum; let \([\alpha,\beta]\) be the largest interval on which \(\phi\) is constant, show that there exists \(\delta>\beta\) close enough to \(\beta\) and such that \(\phi(\delta)=\phi(\beta)\) , by observing that if \(\|u+v\|=\|u-v\|\) , then \(\|u+\gamma^{n}v\|=\|u-\gamma^{n}v\|\) for every rational integer n).
\item
  Show that if \(\|x + y\| = \|x - y\|\) , then \(\|x + \xi y\| = \|x - \xi y\|\) for every real number \(\xi\) . (With the notations of a), observe that \(\phi\) has a relative minimum at the point \(\xi = 0\) , using the fact that \(\phi(\gamma^{n}) = \phi(-\gamma^{n})\) for every rational integer n; deduce that we have \(\phi(\xi) = \phi(-\xi)\) identically, for, otherwise, we get \(\phi(\lambda) = \phi(\mu)\) for two numbers \(\lambda, \mu\) such that \(\lambda + \mu \neq 0\) and that in this case \(\phi\) has a relative minimum at the point \(\frac{1}{2}(\lambda + \mu)\) .
\item
  Deduce from b) that the norm on E is prehilbertian (show first that if \(\|x\| = \|y\|\) , we have \(\|\alpha x + \beta y\| = \|\beta x + 2y\|\) for every pair of real numbers \(\alpha\) , \(\beta\) and that the equality \(\|x + y\| = \|x - y\|\) implies \(\|\alpha x + \beta y\| = \|\alpha x - \beta y\|\) for every pair of real numbers \(\alpha\) , \(\beta\) . Next, show that if \(\|x\| = \|y\| = 1\) and \(\|x + y\| = \|x - y\|\) , we have the relation \(\left\|\left(\alpha^{2} - \beta^{2}\right)x + 2\alpha\beta y\right\| = \alpha^{2} + \beta^{2}\) by using the preceding results, and deduce the conclusion).
\end{enumerate}

\begin{enumerate}
\def\labelenumi{\arabic{enumi})}
\setcounter{enumi}{7}
\tightlist
\item
  Let \(E\) be a real or complex normed space, having the following property: if \(x, y, x', y'\) are four vectors in \(E\) such that
\end{enumerate}

\[
\| x \| = \| x ^ {\prime} \|, \quad \| y \| = \| y ^ {\prime} \|, \quad \| x + y \| = \| x ^ {\prime} + y ^ {\prime} \|,
\]

then \(\| x - y \| = \| x' - y' \|\) . Show that the norm on \(E\) is prehilbertian (use exerc. 7).

\begin{enumerate}
\def\labelenumi{\arabic{enumi})}
\setcounter{enumi}{8}
\item
  Let E be a real hilbertian space, f a continuous linear form on E. Show that on every closed convex subset A of E, the function \(x \mapsto \|x\|^{2} - f(x)\) is bounded below and attains its minimum at a unique point of A.
\item
  Let \(\mathbf{E}\) be a real hilbertian space, \(\mathbf{B}\) a bilinear form on \(\mathrm{E} \times \mathrm{E}\) , \(c_{1}\) , \(c_{2}\) two numbers \(> 0\) such that
\end{enumerate}

\[
\left| \mathrm{B} (x, y) \right| \leqslant c _ {1} \| x \|. \| y \| \quad \text { for   every } x, y \text { in } E;
\]

\[
\left| \mathrm{B} (x, x) \right| \geqslant c _ {2} \| x \| ^ {2} \quad \text {   for   every   } x \in \mathrm{E}.
\]

Show that for every continuous linear form \(f\) on \(\mathbf{E}\) , there exists a unique element \(x_{f} \in \mathbf{E}\) (resp. \(y_{f} \in \mathbf{E}\) ) such that \(f(x) = \mathbf{B}(x_{f}, x)\) (resp. \(f(x) = \mathbf{B}(x, y_{f})\) ) for every \(x \in \mathbf{E}\) .

\begin{enumerate}
\def\labelenumi{\arabic{enumi})}
\setcounter{enumi}{10}
\tightlist
\item
  Let \(\mathbf{E}\) be a hilbertian space, and \((x_{n})\) be a sequence of points of \(\mathbf{E}\) which converges weakly to a point \(a\) . For every \(y \in \mathbf{E}\) , we put
\end{enumerate}

\[
d (y) = \lim _ {n \rightarrow \infty} \inf \| x _ {n} - y \| \quad \text { and } \quad D (y) = \lim _ {n \rightarrow \infty} \sup \| x _ {n} - y \|.
\]

Show that \(d(y)^2 = d(a)^2 + \| y - a\|^2\) and \(\mathrm{D}(y)^2 = \mathrm{D}(a)^2 + \| y - a\|^2\) . If \(\alpha\) and \(\beta\) are two real numbers such that \(0 \leqslant \alpha \leqslant \beta\) , give examples where \(d(a) = \alpha\) and \(\mathrm{D}(a) = \beta\) .

¶ 12) a) Show that there exists a number \(c_{0} > 0\) such that, for every real normed vector space E of dimension \(n\) and every integer \(k \leqslant c_{0} n\) , there exists a hilbertian norm \(x \mapsto \| x \|_{2}\) on E such that \(\| x \|_{2} \leqslant \| x \|\) for all \(x \in \mathbf{E}\) , as well as an orthonormal system \((x_{j})_{1 \leqslant j \leqslant k}\) of \(k\) elements of E (for the hilbertian structure) with norms \(\| x_{j} \| \leqslant 2\) . (Use exerc. 3 of V, p.~60.) b) Let \(n, m\) be two integers \(> 0\) such that \(n \leqslant c_{0} m\) . Let E be a real normed vector space of dimension \(m\) . Show that there exists a vector subspace F of E, of dimension \(n\) , a positive and separating symmetric bilinear form \((x, y) \mapsto \langle x | y \rangle\) on F and an orthonormal basis \(\{a_{1}, a_{2}, ..., a_{n}\}\) of F such that

\[
\frac {1}{2} \sup _ {j} | \langle a _ {j} | x \rangle | \leqslant \| x \| \leqslant \| x \| _ {2}
\]

(where \(\| x\| _2^2 = \langle x|x\rangle\) ) for all \(x\in \mathrm{F}\) . (Apply a) to the dual E' of E.)

\begin{enumerate}
\def\labelenumi{\arabic{enumi})}
\setcounter{enumi}{12}
\tightlist
\item
  \begin{enumerate}
  \def\labelenumii{\alph{enumii})}
  \tightlist
  \item
    Let \((x_{n})_{n\in \mathbb{N}}\) be an infinite sequence in a Banach space E. Show that, in order that the family \((x_{n})\) be summable, it is necessary and sufficient that, for every sequence \((\varepsilon_{n})\) of numbers equal to 1 or to \(-1\) , the series with the general term \((\varepsilon_{n}x_{n})\) is convergent (use GT, III, § 5, exerc. 4).
  \end{enumerate}
\end{enumerate}

\begin{enumerate}
\def\labelenumi{\alph{enumi})}
\setcounter{enumi}{1}
\tightlist
\item
  Let \((x_{j})_{1\leqslant j\leqslant n}\) be a finite sequence of points in a hilbertian space \(E\) show that
\end{enumerate}

\[
2 ^ {- n} \sum_ {(\varepsilon_ {j})} \left(\left\| \sum_ {j = 1} ^ {n} \varepsilon_ {j} x _ {j} \right\| ^ {2}\right) = \sum_ {j = 1} ^ {n} \| x _ {j} \| ^ {2},
\]

where \((\varepsilon_{j})\) ranges over the set of \(2^{n}\) sequences of numbers equal to 1 or to \(-1\) (use the identity of the median, cf.~V, p.~9, formula (14)).

\begin{enumerate}
\def\labelenumi{\alph{enumi})}
\setcounter{enumi}{2}
\tightlist
\item
  Deduce from \(b)\) that if \((x_{i})_{i\in I}\) is a summable family in a hilbertian space \(E\) , the family \((\| x_i\|^2)_{i\in I}\) is summable in \(\mathbf{R}\) .
\end{enumerate}

¶ 14) Let E be an infinite dimensional Banach space. a) Show that for every integer N, there exists a sequence \((b_{j})_{1 \leqslant j \leqslant N}\) of N vectors in E, of norm 1, such that, for every sequence \((\xi_{j})_{1 \leqslant j \leqslant N}\) of N scalars, we have

\[
\left\| \sum_ {j = 1} ^ {\mathrm{N}} \xi_ {j} b _ {j} \right\| ^ {2} \leqslant 4 \sum_ {j = 1} ^ {\mathrm{N}} | \xi_ {j} | ^ {2}
\]

(use exerc. 12, b)).

\begin{enumerate}
\def\labelenumi{\alph{enumi})}
\setcounter{enumi}{1}
\item
  For every sequence \((\lambda_n)_{n\geqslant 1}\) of numbers \(\geqslant 0\) such that \(\sum_{n}\lambda_{n}^{2} < +\infty\) , show that there exists a sequence \((x_{n})_{n\geqslant 1}\) of points of E such that \(\| x_n\| = \lambda_n\) for all \(n\) , and that the series \((x_{n})\) is summable. (Use a) of exerc. 13, a).
\item
  Deduce from b) that in every infinite dimensional Banach space, there exists a commutatively convergent series, that is not absolutely convergent (Dvoretzky-Rogers th.).
\end{enumerate}

\begin{enumerate}
\def\labelenumi{\arabic{enumi})}
\setcounter{enumi}{14}
\tightlist
\item
  Let E be a complex hilbertian space, \(E_{1}\) , \(E_{2}\) two closed vector subspaces of E, \(P_{1}\) , \(P_{2}\) the orthoprojectors from E onto \(E_{1}\) , \(E_{2}\) respectively.
\end{enumerate}

\begin{enumerate}
\def\labelenumi{\alph{enumi})}
\item
  Show that, in order that \(P_{1}\) and \(P_{2}\) commute, it is necessary and sufficient that E is the hilbertian sum of the four subspaces \(E_{1} \cap E_{2}\) , \(E_{1}^{\circ} \cap E_{2}^{\circ}\) , \(E_{1}^{\circ} \cap E_{2}\) , \(E_{1} \cap E_{2}^{\circ}\) (where \(M^{\circ}\) denotes the orthogonal complement of a vector subspace M of E).
\item
  Show that if \(E_{1}\) is finite dimensional and \(\|P_{1}-P_{2}\|<1\) , then \(E_{2}\) has the same dimension as \(E_{1}\) (consider the intersection \(E_{1}^{\circ}\cap E_{2}\) ).
\item
  Show that the endomorphism \(T=(P_{1}-P_{2})^{2}\) of E commutes with \(P_{1}\) and \(P_{2}\) , and that the eigen subspace of T corresponding to the eigenvalue 0 (resp. 1) is the direct sum of the orthogonal subspaces \(E_{1}\cap E_{2}\) and \(E_{1}^{\circ}\cap E_{2}^{\circ}\) (resp. \(E_{1}^{\circ}\cap E_{2}\) and \(E_{1}\cap E_{2}^{\circ}\) ).
\item
  Suppose E is finite dimensional and that \(T=\lambda I\) with \(\lambda\neq0\) . Then \(\lambda>0\) and E is the hilbertian sum of subspaces of dimension \(\leqslant2\) , each of which is stable under \(P_{1}\) and \(P_{2}\) (observe that \(P_{1}-P_{2}\) is hermitian and deduce that E is the hilbertian sum of two subspaces \(E^{+},E^{-}\) such that \(P_{1}.x-P_{2}.x=\sqrt{\lambda}x\) in \(E^{+}\) and \(P_{1}.x-P_{2}.x=-\sqrt{\lambda}x\) in \(E^{-}\) ; then show that for \(x\in E^{+}\) , we have \(P_{1}.x=\frac{1+\sqrt{\lambda}}{2}x+z\) and \(P_{2}.x=\frac{1-\sqrt{\lambda}}{2}x+z\) , with \(z\in E^{-}\) ). * e) Suppose that \(E_{1}\) and \(E_{2}\) are finite dimensional. Show that there exists a family \((\mathrm{F}_{\alpha})_{\alpha\in\mathrm{A}}\) of subspaces of E of dimension \(\leqslant2\) such that E, \(E_{1}\) and \(E_{2}\) are respectively the hilbertian sums of the families \((\mathrm{F}_{\alpha})_{\alpha\in\mathrm{A}},(\mathrm{F}_{\alpha}\cap\mathrm{E}_{1})_{\alpha\in\mathrm{A}}\) and \((\mathrm{F}_{\alpha}\cap\mathrm{E}_{2})_{\alpha\in\mathrm{A}}\) (use c) to reduce to the case when E is finite dimensional, then apply d)). *
\end{enumerate}

\begin{enumerate}
\def\labelenumi{\arabic{enumi})}
\setcounter{enumi}{15}
\tightlist
\item
  Let \(\mathbf{E}\) be a hilbertian space, and \(P\) be a continuous projector on \(\mathbf{E}\) , i.e.~a continuous endomorphism of \(\mathbf{E}\) such that \(P^2 = P\) . Show that for \(P\) to be an orthoprojector, it is necessary and sufficient that \(\| P \| \leqslant 1\) . (To see that the condition is sufficient, consider a vector \(x\) orthogonal to the kernel of \(I - P\) .)
\end{enumerate}

If P has finite rank, show that there exists a closed subspace F of E, with finite codimension, containing \(P(\mathrm{E})\) and such that the restriction of P to F is an orthoprojector.

\begin{enumerate}
\def\labelenumi{\arabic{enumi})}
\setcounter{enumi}{16}
\tightlist
\item
  \begin{enumerate}
  \def\labelenumii{\alph{enumii})}
  \tightlist
  \item
    Let E be a real hilbertian space of dimension 2, \(P_1\) , \(P_2\) two orthoprojectors from E onto the lines \(\mathbf{D}_1\) , \(\mathbf{D}_2\) respectively, assumed distinct. Show that for every \(x \in E\) such that \(\| x \| = 1\) , we have \(\| (P_1 - P_2).x \| = \sin \theta\) , where \(\theta\) is the angle between \(\mathbf{D}_1\) and \(\mathbf{D}_2\) lying between 0 and \(\pi/2\) , and that for every \(y \neq 0\) in E, there exists \(x \neq 0\) such that \((P_1 - P_2).x\) is collinear with \(y\) .
  \end{enumerate}
\end{enumerate}

\begin{enumerate}
\def\labelenumi{\alph{enumi})}
\setcounter{enumi}{1}
\item
  Let E be a real hilbertian space, \(P_{1}\) , \(P_{2}\) two orthoprojectors on E, with respective images \(E_{1}\) , \(E_{2}\) . Show that \(\|P_{1}-P_{2}\|\) is the lower bound of the numbers \(\sin\theta\) , where \(\theta\) is the angle, between 0 and \(\pi/2\) of the two lines \(D_{1}\) , \(D_{2}\) such that \(D_{1}\subset E_{1}\) , \(D_{2}\subset E_{2}\) , \(D_{1}\) and \(D_{2}\) being orthogonal to \(E_{1}\cap E_{2}\) .
\item
  Let \(Q_{1}, Q_{2}\) be two continuous projectors on \(\mathbf{E}\) , with images \(\mathrm{E}_1, \mathrm{E}_2\) , and let \(P_{1}, P_{2}\) be the orthoprojectors onto \(\mathrm{E}_1\) and \(\mathrm{E}_2\) respectively. Show that \(\| P_2 - P_1 \| \leqslant \| Q_2 - Q_1 \|\) . (Observe that \((Q_2 - Q_1)P_2 = (I - Q_1)(P_2 - P_1)\) and use a) and b).
\end{enumerate}

¶ 18) Let E be a real normed space of dimension \(\geqslant 3\) . Suppose that there exists a decreasing bijective mapping \(\omega\) from the set \(\mathfrak{M}\) of closed vector subspaces of E onto itself, such that \(\omega(\omega(M)) = M\) and \(M \cap \omega(M) = \{0\}\) for every \(M \in \mathfrak{M}\) .

\begin{enumerate}
\def\labelenumi{\alph{enumi})}
\item
  Show that there exists a linear mapping \(u\) from \(\mathbf{E}\) onto its dual \(\mathbf{E}'\) defined up to a scalar factor and such that \(u(\mathbf{M}) = (\omega(\mathbf{M}))^\circ\) for all \(\mathbf{M} \in \mathfrak{M}\) . (Considering the case where \(\mathbf{M}\) is 1 dimensional, apply the fundamental th. of projective geometry (A, II, § 9, exerc. 16) by observing that the only automorphism of the field \(\mathbf{R}\) is the identity mapping (GT, IV, § 3, exerc. 3).
\item
  If we put \(\langle x|y\rangle = \langle x,u(y)\rangle\) , show that \(\langle x|x\rangle \neq 0\) for all \(x\neq 0\) and that the relations \(\langle x|y\rangle = 0\) and \(\langle y|x\rangle = 0\) are equivalent. Deduce that \(\langle y|x\rangle = \langle x|y\rangle\) for every pair of points \(x,y\) of E (consider a number \(\lambda \in \mathbf{R}\) such that \(\langle \lambda x + y|x\rangle = 0\) ).
\item
  Show that \(\langle x|x\rangle\) has the same sign on the set of all \(x\neq 0\) ; replacing \(u\) by \(-u\) if necessary, we can then assume that \(\langle x|y\rangle\) is a positive, separating symmetric bilinear form on \(\mathbf{E}\times \mathbf{E}\) . d) Let \(\mathcal{T}_0\) be the initial topology of E. Show that the topology \(\mathcal{T}\) on E, defined by the norm \(\langle x|x\rangle^{1 / 2}\) is finer than the topology \(\mathcal{T}_0\) (observe that the dual of E for \(\mathcal{T}\) contains the dual E' of E for \(\mathcal{T}_0\) ).
\item
  Show that u is a continuous mapping from E onto its dual \(E'\) , for the topologies \(\sigma(E, E')\) and \(\sigma(E', E)\) . Deduce that if E is complete for the initial topology \(T_{0}\) , then u is continuous for \(T_{0}\) and for the strong topology \(\beta(E', E)\) (observe that u transforms every set bounded for \(\sigma(E, E')\) into a bounded set for \(\sigma(E', E)\) ). Deduce that then the topologies T and \(T_{0}\) are identical, and \(\omega(\mathbf{M})\) is the orthogonal complement of \(\mathbf{M}\) for the hilbertian space structure defined on \(\mathbf{E}\) by the form \(\langle x|y\rangle\) (cf.~I, p.~17, th. 1).
\item
  Show that in the space \(\ell^{1}(\mathbf{N})\) , with the norm induced by that of \(\ell^{\infty}(\mathbf{N})\) assigned to it, there exists a bijective mapping \(\mathbf{M}\mapsto\omega(\mathbf{M})\) from \(\mathfrak{M}\) onto itself, having the properties mentioned above (IV, p.~47, exerc. 1).
\end{enumerate}

¶ 19) Let E be an infinite dimensional complex normed space. Suppose that there exists a bijective mapping \(\omega\) from the set \(\mathfrak{M}\) of closed vector subspaces of E onto itself, having the properties listed in exerc. 18.

\begin{enumerate}
\def\labelenumi{\alph{enumi})}
\item
  Show that there exists a semi-linear mapping \(u\) from E onto its dual \(\mathbf{E}'\) (for the automorphism \(\xi \mapsto \overline{\xi}\) of \(\mathbf{C}\) ) defined up to a scalar factor and such that \(u(\mathbf{M}) = (\omega(\mathbf{M}))^\circ\) for every \(\mathbf{M} \in \mathfrak{M}\) . (Proceed as in exerc. 18; using IV, p.~65, exerc. 16, show that \(u\) is a semi-linear mapping relative to the identity automorphism of \(\xi \mapsto \overline{\xi}\) ; finally prove that the first case cannot occur since \(\langle x, u(x) \rangle \neq 0\) for \(x \neq 0\) .)
\item
  If we put \(\langle y|x\rangle = \langle x,u(y)\rangle\) , show that \(\langle x|y\rangle = \overline{\langle y|x\rangle}\) and that \(\langle x|x\rangle\) has the same sign on the set of all \(x\neq 0\) (same method as in exerc. 18).
\item
  Finally show that the topology defined by the norm \(\langle x|x\rangle^{1 / 2}\) is finer than the initial topology \(\mathcal{T}_0\) on E, and that these two topologies are identical when E is complete for \(\mathcal{T}_0\) ; in the latter case, \(\omega (\mathbf{M})\) is the orthogonal complement of M in the hilbertian space E.
\end{enumerate}

\begin{enumerate}
\def\labelenumi{\arabic{enumi})}
\setcounter{enumi}{19}
\tightlist
\item
  Let \(\mathbf{E}\) be a real finite dimensional vector space, and \(\phi\) be a bijective linear mapping from \(\mathbf{E}\) onto its dual \(\mathbf{E}^*\) . Let \(\mathbf{A}\) be a symmetric compact convex set in \(\mathbf{E}\) , with 0 as an interior point; assume that for every \(x\) in the boundary of \(\mathbf{A}\) , the hyperplane with equation \(\langle y - x, \phi(x) \rangle = 0\) is a support hyperplane for \(\mathbf{A}\) .
\end{enumerate}

\begin{enumerate}
\def\labelenumi{\alph{enumi})}
\item
  Let \(f(x) = |\langle x, \phi(x) \rangle|\) , and let \(a\) be a boundary point of \(A\) where \(f(x)\) attains its minimum. Show that for every point \(b\) such that \(\langle b, \phi(a) \rangle = 0\) , we also have \(\langle a, \phi(b) \rangle = 0\) . (Observe that \(\langle x, \phi(x) \rangle \neq 0\) for \(x \neq 0\) , and so we can assume that \(f(x) = \langle x, \phi(x) \rangle \geqslant 0\) ; use the fact that every support hyperplane of \(A\) at the point \(a\) is also a support hyperplane of the set defined by \(f(x) \leqslant f(a)\) .)
\item
  Show that \((x, y) \mapsto \langle x, \phi(y) \rangle\) is a symmetric bilinear form, and that A is identical with the set of all points \(x\) such that \(f(x) \leqslant \gamma\) for a suitable constant \(\gamma\) . (Argue by induction on the dimension of E; with the notations of \(a\) ), consider the hyperplane with the equation \(\langle x, \phi(a) \rangle = 0\) .)
\end{enumerate}

\begin{enumerate}
\def\labelenumi{\arabic{enumi})}
\setcounter{enumi}{20}
\tightlist
\item
  Let E be a finite dimensional complex vector space, and let \(\phi\) be a bijective semi-linear (relative to the automorphism \(\xi \mapsto \overline{\xi}\) of C) mapping from E onto its dual \(E^{*}\) . Let \(\| x\|\) be a norm on E such that, for all \(x \in E\) , we have \(|\langle x, \phi(x) \rangle| = \| x\| \cdot \| \phi(x)\|\) . Show that \((y, x) \mapsto \langle x, \phi(y) \rangle\) is, up to a constant factor, a positive separating hermitian form and that \(\langle x, \phi(x) \rangle = \gamma \| x\|^2\) ( \(\gamma\) constant). (Argue as in exerc. 20.)
\end{enumerate}

¶ 22) Let E be a real normed space of dimension \(\geqslant 3\) , such that, for every homogeneous plane P in E, there exists a continuous projector from E onto P, of norm 1. Show that the norm on E is prehilbertian. With the help of V, p.~60, exerc. 2, reduce to the case where E is of dimension 3, and establish successively the following propositions.

\begin{enumerate}
\def\labelenumi{\alph{enumi})}
\item
  For every homogeneous plane P in E, there exists a unique continuous projector from E onto P, with norm 1, and the kernel of this projection is a homogeneous line D(P) such that P ↦ D(P) is a continuous bijection from the space of homogeneous planes of E into the space of homogeneous lines of E (GT, VI, § 3, No.~5).
\item
  Every point of the sphere \(\mathrm{D}:\| x\| = 1\) in \(\mathrm{E}\) is extremal in the ball \(\mathbf{B}\) of \(\mathrm{E}\) defined by \(\| x\| \leqslant 1\) . (First show that if \(x\in S\) is not extremal, its section \(F_{x}\) in \(\mathbf{B}\) (II, p.~87, exerc. 3) will be 2-dimensional, by considering all the homogeneous planes \(\mathbf{P}\) passing through \(x\) , next prove that this hypothesis is contradictory, by proceeding similarly at a point in \(F_{x}\) where there exists only one support line of \(F_{x}\) in the plane generated by \(F_{x}\) ; the existence of such a point can be established by using II, p.~88, exerc. 7 and p.~88, exerc. 8).
\item
  Every point of the sphere \(\mathbf{S}'\) with equation \(\| x'\| = 1\) in the dual \(\mathrm{E}'\) of \(\mathrm{E}\) is extremal in the ball \(\mathbf{B}'\) of \(\mathrm{E}'\) defined by \(\| x'\| \leqslant 1\) . (First observe that for every homogeneous line \(\mathbf{D}'\) of \(\mathrm{E}'\) , there exists a unique homogeneous plane \(\mathrm{P}'(\mathbf{D}')\) in \(\mathrm{E}'\) such that for every point in \(\mathbf{S}' \cap \mathrm{P}'(\mathbf{D}')\) , the support plane of \(\mathbf{B}'\) at this point (unique by \(a\) ) is parallel to \(\mathbf{D}'\) ; moreover, the mapping \(\mathbf{D}' \mapsto \mathbf{P}'(\mathbf{D}')\) is continuous. Deduce that if \(x' \in \mathbf{S}'\) were not an extremal point in \(\mathbf{B}'\) , its section \(\mathbf{F}_{x'}\) in \(\mathbf{B}'\) would be of dimension at least 2; for this consider all the homogeneous lines \(\mathbf{D}'\) parallel to the support plane of \(\mathbf{B}'\) at the point \(x'\) . Next show that this hypothesis implies a contradiction, by considering a point of strict convexity \(y'\) of \(\mathbf{F}_{x'}\) (II, p.~88, exerc. 8), and the unique homogeneous line \(\mathbf{D}_0'\) parallel to the support line of \(\mathbf{F}_{x'}\) at the point \(y'\) in the plane generated by \(\mathbf{F}_{x'}\) , and prove that the function \(\mathbf{D}' \mapsto \mathbf{P}'(\mathbf{D}')\) would not be continuous for \(\mathbf{D}' = \mathbf{D}_0'\) .) \(d)\) Show that, if three homogeneous planes \(\mathbf{P}_1, \mathbf{P}_2, \mathbf{P}_3\) in E contain the same line \(\Delta\) , then the three lines \(\mathbf{D}(\mathbf{P}_1), \mathbf{D}(\mathbf{P}_2), \mathbf{D}(\mathbf{P}_3)\) are in the same homogeneous plane \(\pi(\Delta)\) (consider the unique support plane of B at a point of intersection of \(\Delta\) and of S). Applying the fundamental theorem of projective geometry (A, II, § 9, exerc. 16) deduce that there exists a bijective linear mapping \(\phi\) from \(\mathbf{E}'\) onto E such that, for all \(x' \in \mathbf{E}'\) , the point \(\phi(x')\) belongs to the line \(\mathbf{D}(\mathbf{P})\) , where \(\mathbf{P}\) is the plane with the equation \(\langle y, x' \rangle = 0\) . Show that for every point \(x' \in \mathbf{S}'\) , the plane with the equation \(\langle \phi(x'), y' - x' \rangle = 0\) is the support plane of \(\mathbf{B}'\) at this point, and conclude by applying V, p.~65, exerc. 20.
\end{enumerate}

\begin{enumerate}
\def\labelenumi{¶ \arabic{enumi})}
\setcounter{enumi}{22}
\item
  Let E be a complex normed vector space of dimension \(\geqslant 3\) , such that for every (complex) homogeneous plane P in E, there exists a continuous projector from E onto P with norm 1. Show that the norm on E is prehilbertian. Using V, p.~60, exerc. 2 reduce to the case where E is of dimension 3 over C, and proceed as in exerc. 22. (For part b) of the proof, consider, for every \(x' \in E'\) such that \(\|x'\| = 1\) , the convex set \(G_{x'}\) of all \(x \in S\) such that \(\langle x, x' \rangle = 1\) , show that if \(G_{x'}\) is not simply 0, it would have dimension at least 3 over R; then in the real affine linear variety generated by \(G_{x'}\) , consider a boundary point of \(G_{x'}\) , where there exists only one support hyperplane (real) of \(G_{x'}\) . Similarly, for part c) of the proof, consider, for all \(x \in S\) , the set \(G_{x}'\) of all \(x' \in S'\) such that \(\langle x, x' \rangle = 1\) and show that \(G_{x}'\) reduces to a point; for this, prove that, if not, the real affine linear variety generated by \(G_{x}'\) will have dimension at least 2 over R, and will contain two linearly independent vectors over C. Conclude using V, p.~65, exerc. 21.)
\item
  In a real normed space \(\mathbf{E}\) of dimension \(\geqslant 3\) , we say that a vector \(y\) is quasi-normal to a vector \(x\) , if for every scalar \(\lambda\) , we have \(\| x + \lambda y \| \geqslant \| x \|\) .
\end{enumerate}

\begin{enumerate}
\def\labelenumi{\alph{enumi})}
\item
  Show that, if the relation « y is quasi-normal to x » is symmetric in x, y, then the norm on E is prehilbertian (show that the condition of V, p.~65, exerc. 22 is satisfied).
\item
  Show that the same conclusion holds, if for every closed homogeneous hyperplane H in E, there exists a vector \(\neq 0\) which is quasi-normal to all the vectors of H. (Same method, apply th. 2 of E, III, § 2, No.~4 to the continuous projectors of norm 1 from the vector subspaces containing P onto a homogeneous plane P, these projections being linearly ordered by the relation of extension.)
\item
  Show that the same conclusion holds if for every vector \(x \neq 0\) in E, there exists a closed hyperplane H such that x is quasi-normal to all the vectors of H. (Reduce to the case where E is of dimension 3, and apply V, p.~65, exerc. 22 to the dual of E).
\item
  Show that the same conclusion holds if, when z is quasi-normal to x and y, then z is quasi-normal to \(x + y\) (apply V, p.~65, exerc. 22).
\end{enumerate}

\begin{enumerate}
\def\labelenumi{\arabic{enumi})}
\setcounter{enumi}{24}
\tightlist
\item
  \begin{enumerate}
  \def\labelenumii{\alph{enumii})}
  \tightlist
  \item
    Let E be a real normed space and \(x' \neq 0\) a vector in the dual E' of E. Show that for every vector \(y\) in the hyperplane \(x'^{-1}(0)\) to be quasi-normal to \(x\) (exerc. 24), it is necessary and sufficient that \(\langle x, x' \rangle = \| x \| \cdot \| x' \|\) .
  \end{enumerate}
\end{enumerate}

\begin{enumerate}
\def\labelenumi{\alph{enumi})}
\setcounter{enumi}{1}
\item
  Deduce from \(a)\) that for all \(x \neq 0\) in \(\mathbf{E}\) , there exists a closed homogeneous hyperplane \(\mathbf{H}\) of \(\mathbf{E}\) such that every vector \(y \in \mathbf{H}\) is quasi-normal to \(x\) .
\item
  If \(x, y\) are two points in \(E\) and \(x \neq 0\) , then there exists a scalar \(\alpha\) such that \(\alpha x + y\) is quasi-normal to \(x\) .
\end{enumerate}

\begin{enumerate}
\def\labelenumi{\arabic{enumi})}
\setcounter{enumi}{25}
\tightlist
\item
  A real normed space E is said to be smooth if all the points of the unit sphere in E are points of smoothness (II, p.~87, exerc. 6) of the unit ball. For this to be so, it is necessary and sufficient that there exists a unique positively homogeneous mapping f from E - \{0\} into E' - \{0\}, such that \(\|f(x)\| = 1\) for \(\|x\| = 1\) , and that \(\langle x, f(x) \rangle = \|x\| \cdot \|f(x)\|\) . Show that the following properties are equivalent :
\end{enumerate}

α) E is smooth.

\(\beta)\) For all \(x \neq 0\) in E and all \(y \in \mathbf{E}\) , there exists a unique scalar \(\alpha\) such that \(\alpha x + y\) is quasi-normal to \(x\) .

\(\gamma)\) For every \(x \in E\) , if \(y\) and \(z\) are quasi-normal to \(x\) , \(y + z\) is quasi-normal to \(x\) .

(To see that \(\gamma\) ) implies \(\beta\) ), observe that if \(\alpha x + y\) and \(\beta x + y\) are quasi-normal to \(x\) , then \((\alpha - \beta)x\) is quasi-normal to \(x\) .)

\begin{enumerate}
\def\labelenumi{\arabic{enumi})}
\setcounter{enumi}{26}
\item
  A real normed space E is said to be strictly convex if all the points of the unit sphere are points of strict convexity (II, p.~87, exerc. 6) of the unit ball. Show that, for E to be strictly convex, it is necessary and sufficient that for all \(x \neq 0\) in E and for all \(y \in E\) , there exists a unique scalar \(\alpha\) such that \(x\) is quasi-normal to \(x + y\) . (Observe that the mapping \(t \mapsto \| tx + y \|\) is convex in R.)
\item
  Let \(\mathbf{E}\) be a normed space, \(\mathbf{E}'\) its dual.
\end{enumerate}

\begin{enumerate}
\def\labelenumi{\alph{enumi})}
\tightlist
\item
  Show that if \(\mathrm{E}'\) is smooth (V, p.~66, exerc. 26), E is strictly convex (if \(x\) and \(y\) are such that \(x \neq y\) , \(\| x \| = \| y \| = \| \frac{1}{2}(x + y) \| = 1\) , consider an \(x' \in \mathrm{E}'\) such that \(\| x' \| = 1\) and
\end{enumerate}

\[
\left\langle \frac {1}{2} (x + y), x ^ {\prime} \right\rangle = 1).
\]

\begin{enumerate}
\def\labelenumi{\alph{enumi})}
\setcounter{enumi}{1}
\tightlist
\item
  Show that if \(\mathbf{E}'\) is strictly convex, \(\mathbf{E}\) is smooth.
\end{enumerate}

¶ 29) Let E be a normed space, E' its dual. A mapping f from E - \{0\} into E' - \{0\} is said to be a support mapping if it is positively homogeneous, and if for every x ∈ E such that \textbar\textbar x\textbar\textbar{} = 1, we have \textbar\textbar f(x)\textbar\textbar{} = 1 and ⟨x, f(x)⟩ = 1. For E to be smooth (V, p.~66, exerc. 26), it is necessary and sufficient that there exists a unique support mapping from E - \{0\} into E' - \{0\}.

Let \(S\) be the unit sphere in \(E\) , \(S'\) the unit sphere in \(E'\) , and let \(x_0 \in S\) . The following conditions are equivalent:

α) \(x_{0}\) is a point of smoothness of the unit ball in E,

β) There exists a support mapping \(f\) whose restriction to \(S\) is continuous at the point \(x_0\) when \(S\) is assigned the norm topology, and \(S'\) the weak topology \(\sigma(E', E)\) .

\(\gamma)\) For every \(y \in \mathbf{E}\) , the mapping \(t \mapsto \| x_0 + ty\|\) has a derivative at the point \(t = 0\) .

(To see that \(\alpha\) ) implies \(\beta\) ), argue by reductio ad absurdum, using the weak compactness of the unit ball in E'. To see that \(\beta\) ) implies \(\gamma\) ), reduce to the case where E is 2-dimensional and use the fact that \(t \mapsto \|x_0 + ty\|\) is convex.)

Then every support mapping is continuous at the point \(x_0\) .

\begin{enumerate}
\def\labelenumi{\arabic{enumi})}
\setcounter{enumi}{29}
\tightlist
\item
  Let E be a Banach space, \(E'\) its strong dual, \(E''\) the strong dual of \(E'\) , \(E'''\) the strong dual of \(E''\) , \(E^{IV}\) the strong dual of \(E'''\) .
\end{enumerate}

\begin{enumerate}
\def\labelenumi{\alph{enumi})}
\item
  Suppose that E is non-reflexive; then there exists \(x' \in E'\) such that \(\| x' \| = 1\) , but that we do not have \(\langle x, x' \rangle = 1\) for any \(x \in E\) with \(\| x \| = 1\) (IV, p.~57, exerc. 25). On the other hand, there exists a sequence \((x_n')\) of points of \(E'\) such that \(\| x_n' \| = 1\) , tending strongly to \(x'\) , and a sequence \((x_n)\) of points of E such that \(\| x_n \| = 1\) and \(\langle x_n, x_n' \rangle = 1\) for all \(n\) (II, p.~77, exerc. 4). Show that in \(E''\) , the sequence \((x_n)\) does not converge to any point for the topology \(\sigma(E'', E''')\) (observe that otherwise it would converge to a point \(x \in E\) for which \(\langle x, x' \rangle = 1\) ).
\item
  Show that it is not possible that \(x'\) and \(x_n'\) are points of smoothness of the unit sphere in \(\mathrm{E}'''\) (observe that \(x_n\) , considered as an element of \(\mathrm{E}^{\mathrm{IV}}\) will be the unique element \(x_n^{\mathrm{IV}} \in \mathrm{E}^{\mathrm{IV}}\) such that \(\| x_n^{\mathrm{IV}} \| = 1\) and \(\langle x_n', x_n^{\mathrm{IV}} \rangle = 1\) and use exerc. 29).
\item
  Conclude that if \(\mathrm{E}'''\) is smooth, or if \(\mathrm{E}^{\mathrm{IV}}\) is strictly convex, then \(\mathrm{E}\) is necessarily reflexive.
\end{enumerate}

\begin{enumerate}
\def\labelenumi{\arabic{enumi})}
\setcounter{enumi}{30}
\tightlist
\item
  A normed space E (real or complex) is said to be uniformly convex if, for every \(\varepsilon\) such that \(0 < \varepsilon < 2\) , there exists \(\delta > 0\) such that the relations \(\|x\| \leqslant 1\) , \(\|y\| \leqslant 1\) , \(\|x - y\| \geqslant \varepsilon\) in E imply \(\left\|\frac{1}{2}(x + y)\right\| \leqslant 1 - \delta\) . A uniformly convex space is strictly convex (V, p.~67, exerc. 27). We say that E is uniformly smooth if, for every \(\varepsilon > 0\) , there exists \(\eta > 0\) such that the relations \(\|x\| \geqslant 1\) , \(\|y\| \geqslant 1\) , \(\|x - y\| \leqslant \eta\) imply the inequality \(\|x + y\| \geqslant \|x\| + \|y\| - \varepsilon \|x - y\|\) . This is equivalent to: for every \(\varepsilon > 0\) , there exists \(\rho > 0\) such that the relations \(\|x\| = 1\) , \(\|y\| \leqslant \rho\) imply the inequality
\end{enumerate}

\[
\| x + y \| + \| x - y \| \leqslant 2 + \varepsilon \| y \|.
\]

A uniformly smooth space is smooth (V, p.~66, exerc. 26).

\begin{enumerate}
\def\labelenumi{\alph{enumi})}
\item
  Show that if E is uniformly convex, its strong dual \(E'\) is uniformly smooth, and that if E is uniformly smooth, then \(E'\) is uniformly convex; the restriction of the unique support mapping (V, p.~67, exerc. 29) to the unit sphere S of E is a mapping from S into the unit sphere \(S'\) of \(E'\) , which is continuous for the norm topologies of E and \(E'\) .
\item
  Show that if E is uniformly convex, and, if a filter \(\mathfrak{F}\) on E converges to \(x_{0}\) for the topology \(\sigma(\mathrm{E}, \mathrm{E}')\) and is such that \(\lim_{\mathfrak{F}} \|x\| = \|x_{0}\|\) , then \(\mathfrak{F}\) converges to \(x_{0}\) for the initial topology of E.
\item
  Show that a Banach space which is uniformly convex or uniformly smooth is reflexive (use b) and c) and also IV, p.~60, exerc. 12). (cf.~V, p.~72, exerc. 14.)
\item
  Generalize the first part of th. 1 of V, p.~10, and also cor. 1 and 2 of V, p.~11 to uniformly convex Banach spaces.
\end{enumerate}

\begin{enumerate}
\def\labelenumi{\arabic{enumi})}
\setcounter{enumi}{31}
\tightlist
\item
  Let \(\mathrm{E}\) be a normed space (real or complex) of dimension \(\geqslant 2\) , such that, for every \(\varepsilon\) such that \(0 < \varepsilon < 2\) , the relations \(\| x \| = 1\) , \(\| y \| = 1\) , \(\| x - y \| \geqslant \varepsilon\) in \(\mathrm{E}\) imply the inequality \(\| \frac{1}{2}(x + y) \| \leqslant \left(1 - \frac{\varepsilon^2}{4}\right)^{1/2}\) . Show that the norm on \(\mathrm{E}\) is prehilbertian. (Reduce to the case when E is real and 2-dimensional, and argue as in V, p.~61; exerc. 4, a).
\end{enumerate}

¶ 33) Let E be a uniformly convex Banach space (V, p.~67, exerc. 31). Then there exists a number θ such that \(\frac{3}{4} \leqslant \theta < 1\) and such that the relation \(\|x - y\| \geqslant \frac{1}{2} \sup (\|x\|, \|y\|)\) in E implies \(\left\|\frac{1}{2}(x + y)\right\| \leqslant \theta \sup (\|x\|, \|y\|)\) .

\begin{enumerate}
\def\labelenumi{\alph{enumi})}
\tightlist
\item
  Let \((x_{n})\) be a sequence of points of E such that \(\| x_{n}\| \leqslant \mathbf{M}\) and such that the sequence tends to 0 for \(\sigma (\mathrm{E},\mathrm{E}^{\prime})\) . Show that, if for an index \(p\) , \(\| x_p\| \geqslant \frac{1}{2}\mathbf{M}\) , then there exists \(q > p\) such that \(\| x_p - x_q\| >\frac{1}{2}\mathbf{M}\) , and consequently \(\left\| \frac{1}{2} (x_p + x_q)\right\| \leqslant \theta \mathbf{M}\) (argue by reductio ad absurdum, by observing that for all \(x^{\prime}\in \mathrm{E}^{\prime}\) such that \(\| x^{\prime}\| = 1\) , we have \(\langle x_p,x'\rangle = \lim_{n\to \infty}\langle x_p - x_n,x'\rangle\) ).
\end{enumerate}

Deduce that there exists a strictly increasing mapping \(\ell\) from \(\mathbf{N}\) into itself such that \(\left\| \frac{1}{2} (x_{\ell(2n)} + x_{\ell(2n+1)})\right\| \leqslant \mathrm{M}\theta\) , and such that if \(x_{n}^{(1)} = \frac{1}{2} (x_{\ell(2n)} + x_{\ell(2n+1)})\) , the sequence \((x_{n}^{(1)})\) tends to 0 for \(\sigma(\mathrm{E},\mathrm{E}')\) and that \(\| x_n^{(1)}\| \leqslant \mathrm{M}\theta\) for all \(n\) .\\
b) Show that there exists a sequence \((x_{n_k})\) extracted from \((x_{n})\) such that, if we put \(y(k) = x_{n_k}\) , we have the following property: for every integer \(p > 1\) , every integer \(q < p\) and every integer \(i\) such that \(1 \leqslant i \leqslant 2^{p - q}\) ,

\[
\left\| y ((i - 1) 2 ^ {q} + 1) + y ((i - 1) 2 ^ {q} + 2) + \dots + y (i 2 ^ {q}) \right\| \leqslant \mathrm{M} \theta^ {q}.
\]

(Iterate the procedure of \(a\) ) by constructing a sequence \((x_{n}^{(k + 1)})\) from the sequence \((x_{n}^{(k)})\) in the same way as \((x_{n}^{(1)})\) is constructed from \((x_{n})\) ; then use a suitable « method of diagonalization ».) \(c)\) Let \(r\) and \(q\) be two integers \(>1\) . Deduce from \(b\) ) that if \(r2^{q}\leqslant k\leqslant (r + 1)2^{q}\) , then

\[
\left\| x _ {n _ {1}} + x _ {n _ {2}} + \dots + x _ {n _ {k}} \right\| \leqslant (2 ^ {q} - 1) \mathrm{M} + 2 ^ {q} \mathrm{M} + (r - 1) 2 ^ {q} \mathrm{M} \theta^ {q}
\]

(decompose the sum on the left into several subsets, by varying h from 1 to \(2^{q}\) , then from \((j - 1)\) \(2^{q} + 1\) to \(j2^{q}\) for \(2 \leqslant j \leqslant r\) , then from \(r2^{q} + 1\) to k).

\begin{enumerate}
\def\labelenumi{\alph{enumi})}
\setcounter{enumi}{3}
\tightlist
\item
  Show that for every bounded sequence \((x_{n})\) in E, there exists an extracted sequence \((x_{n_{k}})\) such that the sequence of the averages \((x_{n_{1}} + \cdots + x_{n_{k}})/k\) converges for the initial topology of E (the Banach-Saks-Kakutani theorem). (Using the fact that E is reflexive, reduce to the case when the sequence \((x_{n})\) converges to 0 for \(\sigma(\mathrm{E}, \mathrm{E}')\) , and use c) for q and r large enough.)
\end{enumerate}

\begin{enumerate}
\def\labelenumi{\arabic{enumi})}
\setcounter{enumi}{33}
\tightlist
\item
  Let E be a Banach space and K be a convex, bounded subset that is closed for \(\sigma(\mathrm{E},\mathrm{E}')\) . Suppose that for every sequence \((x_{n})\) in K, there exists an extracted sequence \((x_{n_{k}})\) such that the sequence of averages \((x_{n_{1}}+\cdots+x_{n_{k}})/k\) is convergent for \(\sigma(\mathrm{E},\mathrm{E}')\) . Show that for every continuous linear form \(x'\) on E, there exists an element x of K such that \(\langle x,x'\rangle=\sup_{y\in K} \langle y,x'\rangle\)
\end{enumerate}

(apply the hypothesis to a sequence \((x_{n})\) of points of \(\mathbf{K}\) such that \(\langle x_n, x'\rangle\) tends to \(\sup_{y\in \mathbf{K}}\langle y, x'\rangle\) ).

Deduce that if \(\mathbf{E}\) has the property of exerc. 33, \(d\) ), then \(\mathbf{E}\) is a reflexive space (cf.~IV, p.~57, exerc. 25).

* 35) Let E denote a real, finite dimensional hilbertian space of dimension \(n\) , S the unit sphere of E and \(m\) the unique positive measure of norm 1 on S which is invariant under the group of automorphisms of E. Consider S as a metric space in which the distance is defined by \(d(x, y) = \operatorname{Arc} \cos \langle x | y \rangle\) . For every \(x \in S\) and \(r \geqslant 0\) , let \(\mathbf{B}(x, r)\) denote the set of all points y in S such that \(d(x, y) \leqslant r\) ; for every subset A of S and every real number \(r \geqslant 0\) , let \(A_r\) be the set of all points x in S such that \(d(x, A) \leqslant r\) .

\begin{enumerate}
\def\labelenumi{\alph{enumi})}
\item
  Given two closed subsets A and B of S, let \(\delta(A, B)\) be the lower bound of the set of all real numbers \(r \geqslant 0\) such that \(A \subset B_r\) and \(B \subset A_r\) . Show that \(\delta\) is a distance on the set \(\mathcal{F}\) of all closed subsets of S, and that \(\mathcal{F}\) is a compact metric space for this distance. Show that the mapping \(A \mapsto m(A)\) from \(\mathcal{F}\) into \(\mathbf{R}\) is upper semi-continuous.
\item
  Let \(x_0\) be a point in \(S\) , let \(H\) be the hyperplane in \(E\) orthogonal to \(x_0\) , \(x_1\) a point in \(H\) and \(\gamma\) the arc of the circle joining \(x_0\) to \(-x_0\) , passing through \(x_1\) , i.e.~the set of all points in \(S\) of the form \(x_0 \sin \theta + x_1 \cos \theta\) with \(|\theta| \leqslant \pi / 2\) . For every \(y \in \gamma\) , let \(H_y = H + y\) and \(S_y = S \cap H_y\) ; let \(m_y\) denote the unique positive measure of norm 1 on \(S_y\) which is invariant under the group of all automorphisms of \(E\) which leave \(x_0\) fixed.
\end{enumerate}

Let A be a closed subset of S and \(\gamma'\) the set of all points in \(\gamma\) such that \(A \cap S_{p}\) is non-empty. For every \(y \in \gamma'\) , there exists a unique real number \(r(y)\) such that \(0 \leqslant r(y) \leqslant \pi\) and such that \(m_{y}(A \cap S_{y}) = m_{y}(B(y, r(y)) \cap S_{y})\) ; let \(s_{\gamma}(A)\) be the union of the sets \(B(y, r(y)) \cap S_{y}\) as y ranges over \(\gamma'\) . Prove that \(s_{\gamma}(A)\) is closed and that \(m(A) = m(s_{\gamma}(A))\) .

\begin{enumerate}
\def\labelenumi{\alph{enumi})}
\setcounter{enumi}{2}
\tightlist
\item
  For every closed subset A of S, the infimum \(r(\mathbf{A})\) of the set of all real numbers \(r \geqslant 0\) for which there exists an \(x \in S\) with \(\mathbf{A} \subset \mathbf{B}(x, r)\) is called the radius of A. Let \(\mathbf{M}(\mathbf{A})\) denote the set of all closed subsets C of S such that \(m(\mathbf{C}) = m(\mathbf{A})\) and \(m(\mathbf{C}_{\varepsilon}) \leqslant m(\mathbf{A}_{\varepsilon})\) for every \(\varepsilon > 0\) . Show that the following conditions are equivalent for every pair (A, B) of closed subsets of S :
\end{enumerate}

\begin{enumerate}
\def\labelenumi{(\roman{enumi})}
\item
  \(m(\mathbf{A}) = m(\mathbf{B})\) and \(\mathbf{B}\) is of the form \(\mathbf{B}(x,r)\) with \(x\in \mathbf{S}\) and \(r\geqslant 0\)
\item
  \(\mathbf{B} \in \mathbf{M}(\mathbf{A})\) and \(r(\mathbf{B}) \leqslant r(\mathbf{C})\) for every subset \(\mathbf{C}\) of \(\mathbf{A}\) belonging to \(\mathbf{M}(\mathbf{A})\) . (Arguing by induction on \(n\) , we deduce from \(b\) ) that \(s_{\gamma}(\mathbf{A})\) belongs to \(\mathbf{M}(\mathbf{A})\) for every closed subset \(\mathbf{A}\) of \(\mathbf{S}\) ; if \(r > 0\) is such that \(\mathbf{A} \subset \mathbf{B}(x_1, r)\) show that every point of the boundary of \(\mathbf{B}(x_1, r)\) in \(\mathbf{S}\) which belongs to \(s_{\gamma}(\mathbf{A})\) also belongs to \(\mathbf{A}\) .)
\end{enumerate}

* 36) The notations are the same as in exerc. 35.

\begin{enumerate}
\def\labelenumi{\alph{enumi})}
\item
  Let \(a\) be a vector of norm 1 in \(\mathbf{E}\) , \(\mathrm{K}_{\varepsilon}\) the set of all \(x \in \mathbf{S}\) such that \(|\langle x|a\rangle| \leqslant \sin \varepsilon\) and \(\mathrm{L}_{\varepsilon}\) the set of all \(x \in \mathbf{S}\) such that \(d(x, \mathrm{S}_a) \geqslant \varepsilon\) (where \(\mathrm{S}_a\) is the set of all points of \(\mathbf{S}\) orthogonal to \(a\) ). Show that for \(\varepsilon > 0\) small enough, we have \(m(\mathrm{K}_{\varepsilon}) \geqslant 4e^{-n\varepsilon^2 / 2}\) and \(m(\mathrm{L}_{\varepsilon}) \leqslant 4e^{-n\varepsilon^2 / 2}\) (we observe that the image of the measure \(m\) under the mapping \(x \mapsto \langle x|a\rangle\) from \(\mathbf{S}\) into the interval \([-1, 1]\) of \(\mathbf{R}\) is of the form \(c_n(1 - t^2)^{(n-3)/2}dt\) with a suitable constant \(c_n > 0\) ).
\item
  Let f be a continuous mapping from S into R and \(\mathbf{M}(f)\) a real number such that the set of all \(x \in S\) for which \(f(x) \leqslant \mathbf{M}(f)\) (resp. \(f(x) \geqslant \mathbf{M}(f)\) ) has a measure \(\geqslant \frac{1}{2}\) for m. Let B be the set of all \(x \in S\) such that \(f(x) = \mathbf{M}(f)\) . Deduce from a) that, for every \(\varepsilon > 0\) small enough, the set of all \(x \in S\) such that \(d(x, B) \geqslant \varepsilon\) , has a measure for m at most equal to \(4e^{-n\varepsilon^{2}/2}\) .
\item
  For every \(\varepsilon > 0\) , let \(h(n, \varepsilon)\) be the smallest integer \(h \geqslant 1\) for which there exist points \(x_1, \ldots, x_h\) in S such that \(S = \bigcup_{i=1}^{h} B(x_i, \varepsilon)\) show that \(\lim_{\varepsilon \to 0} (\log h(n, \varepsilon)) / |\log \varepsilon| = n\) .
\item
  Recall that E is a real hilbertian space of dimension n.~Let k be a positive integer and \(\varepsilon\) , \(\varepsilon'\) two strictly positive numbers such that \(4h(k, \varepsilon) < e^{n\varepsilon'^{2}/2}\) . Let f be a mapping from S into R such that \(|f(x) - f(y)| \leqslant \|x - y\|\) for all x, y in S, and \(\mathbf{M}(f)\) a real number satisfying the relation stated in b). Show that there exists a vector subspace F of E, of dimension k, satisfying the following condition: for every \(x \in F \cap S\) , there exists a point y in \(F \cap S\) such that \(\|x - y\| < \varepsilon\) and \(|f(y) - \mathbf{M}(f)| \leqslant \varepsilon'\) .
\end{enumerate}

* 37) Let \(\mathrm{E}\) be a real hilbertian space of dimension \(n\) and let \(\gamma\) be a positive measure on \(\mathrm{E}\) such that \(\int_{\mathrm{E}} e^{i\langle x|y\rangle} d\gamma(y) = \exp(-\pi \|x\|^2/2)\) for all \(x \in \mathrm{E}\) (INT, IX, § 6, No.~5). Let \(m\) be the unique positive measure of norm 1 on the unit sphere S of E which is invariant under the group of automorphisms of E.

\begin{enumerate}
\def\labelenumi{\alph{enumi})}
\item
  Let \(p\) be a continuous function on \(\mathbf{E}\) , satisfying \(p(t,x) = t.p(x)\) for all \(x\in S\) and all real positive \(t\) . Show that \(\int_{\mathbf{E}}pd\gamma = c_n\int_{\mathbf{S}}pdm\) with \(c_{n} = \pi^{1 / 2}\Gamma (n / 2) / \Gamma ((n + 1) / 2)\) .
\item
  Show that there exists a constant \(\mathbf{C} > 0\) , independent of \(n\) , such that
\end{enumerate}

\[
\int_ {E} \sup _ {1 \leqslant i \leqslant n} | \langle x | e _ {i} \rangle | d \gamma (x) \geqslant C. (\log n) ^ {1 / 2}
\]

for every orthonormal basis \((e_1, \dots, e_n)\) of E.

\begin{enumerate}
\def\labelenumi{\alph{enumi})}
\setcounter{enumi}{2}
\tightlist
\item
  Let \(\varepsilon > 0\) and let \(k\) be a positive integer. Deduce from \(b\) ) and exercises 12 (V, p.~63) and 36, \(d\) ) that if \(n\) is large enough, then for even real normed space V of dimension \(n\) , there exists a vector subspace W of V, of dimension \(k\) , satisfying the following property: there exists a real hilbertian space \(\mathbf{W}_1\) , of dimension \(k\) , and a bijective linear mapping \(u\) from W onto \(\mathbf{W}_1\) such that \(\sup (\| u \|, \| u^{-1} \|) \leqslant 1 + \varepsilon\) .
\end{enumerate}

